\section{Вывод}

В ходе работы было изучено устройство и основные принципы действия отражательного клистрона. Для математического описания использовалась теория тороидального резонатора, входящего в состав клистрона. Рассматривался процесс динамического управления плотностью в электронном пучке, возбуждение переменным конвекционным током пучка колебаний в резонаторе. Рабочий режим соответствовал следующим напряжениям:
\begin{itemize}
	\item \(U_{\text{рез}} = 90\) В
	\item \(U_{\text{уск}} = 126\) В
	\item \(U_{\text{отр}} = 42\) В
\end{itemize}

Получены зависимости тока в цепи детектора от напряжения на отражателе при постоянном напряжении на резонаторе и от напряжения на резонаторе при постоянном напряжении на отражателе. При помощи волномера измерена длина волны колебаний, создаваемых клистроном. Также снята зависимость тока в цепи детектора от тока электронного пучка для трёх зон генерации.